\section{Diskussion}
\label{sec:Discussion}
Die gemessenen Waveforms zeigen das qualitativ erwartete Verhalten: 
Mit steigender Bias-Spannung nehmen Signaldauer ab und Signalstärke zu. 
Der Vergleich zwischen Vorder- und Rückseite bestätigt, dass bei der Vorderseite hauptsächlich Elektronen das lange Signal erzeugen, 
während es bei der Rückseite die Löcher sind.

Bei niedrigen Spannungen (z.B. $|V| \lesssim \SI{60}{\volt}$) ist die Signalform stark verrauscht und zeigt teilweise nicht-monotone Verläufe,
was bei der bestimmung der Peaks und der Integration über die Signaldauer zu Unsicherheiten führt.
Dieser effekt ist besonders bei der Rückseite (Löcher) ausgeprägt, was die größere Unsicherheit bei der Löchermobilität erklären könnte.
Hier könnte auch die Erklärung für die ersten Messpunkte bei der Signaldauer und Ladung liegen, 
die teilweise stark von der theoretischen Kurve abweichen. 
Die bestimmten Mobilitäten
\begin{align*}
  \mu_e &= \SI{1373 \pm 15}{\centi\meter\squared\per\volt\per\second}\,, \\
  \mu_h &= \SI{381 \pm 3}{\centi\meter\squared\per\volt\per\second}
\end{align*}
stimmen gut mit den Literaturwerten für intrinsisches Silizium bei Raumtemperatur überein. 
Der Literaturwert für Elektronen beträgt $\mu_e^\text{lit} \approx \SI{1400}{\centi\meter\squared\per\volt\per\second}$ \cite{Wermes}, 
was einer relativen Abweichung von ca. $\SI{1.7}{\percent}$ entspricht. 
Die gemessene Elektronenmobilität liegt damit im erwarteten Bereich und validiert die verwendete Messmethode und Auswertung.
Die Löchermobilität $\mu_h = \SI{381 \pm 3}{\centi\meter\squared\per\volt\per\second}$ weicht etwas stärker vom typischen Literaturwert $\mu_h^\text{lit} \approx \SI{450}{\centi\meter\squared\per\volt\per\second}$ ab, 
was einer relativen Abweichung von ca. $\SI{15}{\percent}$ entspricht. 

Weitere mögliche Fehlerquellen könnten beim Laser liegen, z.B. eine ungleichmäßige Intensitätsverteilung oder eine ungenaue Kalibrierung der Laserintensität.
Zusätzlich könnte ein zu schwacher Laser für die Rückseite zu einem schlechten Signal-Rausch-Verhältnis führen, 
was die Unsicherheit bei der Bestimmung der Signaldauer und Ladung erhöht und eine Auwertung der Rückseite erschwert.